% Ad Atticum 12.1 (ad-atticum-12-01)
% Source: https://raw.githubusercontent.com/PerseusDL/canonical-latinLit/master/data/phi0474/phi057/phi0474.phi057.perseus-lat1.xml
% Fetched by scripts/fetch_latin.py

% Dateline (Perseus): Scr. in Arpinati viii K. Dec. a. 708 (46).

\ciceroLetterOpener{CICERO ATTICO salutem}

\ciceroSection{1} undecimo die, postquam a te discesseram, hoc litterularum exaravi egrediens e villa ante lucem atque eo die cogitabam in Anagnino, postero autem in Tusculano, ibi unum diem; v Kalend. igitur ad constitutum. atque utinam continuo ad complexum meae Tulliae, ad osculum Atticae possim currere! quod quidem ipsum scribe, quaeso, ad me ut, dum consisto in Tusculano, sciam quid garriat, sin rusticatur, quid scribat ad te; eique interea aut scribes salutem aut nuntiabis itemque Piliae. et tamen etsi continuo congressuri sumus, scribes ad me si quid habebis.

\ciceroSection{2} cum complicarem hanc epistulam, noctuabundus ad me venit cum epistula tua tabellarius; qua lecta de Atticae febricula scilicet valde dolui. reliqua quae exspectabam ex tuis litteris cognovi omnia; sed quod scribis igniculum matutinum γεροντικόν , γεροντικώτερον est memoriola vacillare. ego enim iiii Kal. Axio dederam, tibi iii, Quinto quo die venissem, id est v Kal. Hoc igitur habebis, novi nihil. quid ergo opus erat epistula? quid , cum coram sumus et garrimus quicquid in buccam? est profecto quiddam λέσχη , quae habet, etiam si nihil subest, conlocutione ipsa suavitatem.
